\documentclass{csexam}

\usepackage{listings}
\usepackage{xcolor}

% Python code styling
\lstset{
    language=Python,
    basicstyle=\ttfamily\small,
    keywordstyle=\color{blue}\bfseries,
    stringstyle=\color{red},
    commentstyle=\color{green!50!black}\itshape,
    showstringspaces=false,
    breaklines=true,
    frame=single,
    numbers=left,
    numberstyle=\tiny\color{gray}
}

% Course information
\course{CSCI 128: Coding for Problem Solving}
\courseshort{CSCI 128}
\exam{Quiz 5}
\examdate{Winter 2026}

\begin{document}

\maketitle

\begin{rules}
\textbf{Instructions:}
\begin{itemize}
    \item This quiz covers material from Lab 4 (Functions and Image Transformations)
    \item Answer all questions in the space provided
    \item Show your work where applicable
    \item Time limit: 15 minutes
\end{itemize}
\end{rules}

\vspace{0.2in}

\showgrade

\newpage

\begin{questions}

\question[2] Write the basic syntax for defining a function in Python that takes two parameters and returns their sum.
\vspace{1.5in}

\question[3] What is the transformation formula for horizontally mirroring an image? If a pixel is at position $(x, y)$ in an image with width $W$ and height $H$, what is its new position after horizontal mirroring?

\vspace{0.3cm}
New position: \underline{\hspace{3in}}
\vspace{0.8in}

\question[3] When rotating an image 90 degrees clockwise, what happens to the dimensions? If the original image is 640×480, what are the dimensions of the rotated image?

Original dimensions: 640×480 \hspace{0.5in} New dimensions: \underline{\hspace{2in}}

\vspace{0.3cm}
Why do the dimensions change? \underline{\hspace{4in}}
\vspace{0.5in}

\question[2] Given an image with dimensions 800×600, you want to crop a 200×200 region from the center. What should the starting x and y coordinates be?

\vspace{0.3cm}
start\_x = \underline{\hspace{2in}} \hspace{0.5in} start\_y = \underline{\hspace{2in}}

\newpage

\question[3] What does the following function do? Describe the transformation it applies.
\begin{lstlisting}
def mystery_transform(img):
    width, height = img.size
    result = Image.new("RGB", (width, height))
    src_pixels = img.load()
    dst_pixels = result.load()

    for x in range(width):
        for y in range(height):
            new_x = width - 1 - x
            new_y = height - 1 - y
            dst_pixels[new_x, new_y] = src_pixels[x, y]
    return result
\end{lstlisting}

\vspace{1.2in}

\question[2] Why does the order matter when combining transformations? Give an example where crop→rotate produces a different result than rotate→crop.

\vspace{1.5in}

\end{questions}

\end{document}
